\hypertarget{host-management-with-openbmc}{%
\section{Host Management with
OpenBMC}\label{host-management-with-openbmc}}

This document describes the host-management interfaces of the OpenBMC
object structure, accessible over REST.

Note: Authentication

See the details on authentication at
\href{https://github.com/openbmc/docs/blob/master/REST-cheatsheet.md\#establish-rest-connection-session}{REST-cheatsheet}.

This document uses token based authentication method:

\begin{verbatim}
export bmc=xx.xx.xx.xx
export token=`curl -k -H "Content-Type: application/json" -X POST https://${bmc}/login -d '{"username" :  "root", "password" :  "0penBmc"}' | grep token | awk '{print $2;}' | tr -d '"'`
curl -k -H "X-Auth-Token: $token" https://${bmc}/xyz/openbmc_project/...
\end{verbatim}

\hypertarget{inventory}{%
\subsection{Inventory}\label{inventory}}

The system inventory structure is under the
\texttt{/xyz/openbmc\_project/inventory} hierarchy.

In OpenBMC the inventory is represented as a path which is hierarchical
to the physical system topology. Items in the inventory are referred to
as inventory items and are not necessarily FRUs (field-replaceable
units). If the system contains one chassis, a motherboard, and a CPU on
the motherboard, then the path to that inventory item would be:

\texttt{inventory/system/chassis0/motherboard0/cpu0}

The properties associated with an inventory item are specific to that
item. Some common properties are:

\begin{itemize}
\tightlist
\item
  \texttt{Version}: A code version associated with this item.
\item
  \texttt{Present}: Indicates whether this item is present in the system
  (True/False).
\item
  \texttt{Functional}: Indicates whether this item is functioning in the
  system (True/False).
\end{itemize}

The usual \texttt{list} and \texttt{enumerate} REST queries allow the
system inventory structure to be accessed. For example, to enumerate all
inventory items and their properties:

\begin{verbatim}
curl -k -H "X-Auth-Token: $token" https://${bmc}/xyz/openbmc_project/inventory/enumerate
\end{verbatim}

To list the properties of one item:

\begin{verbatim}
curl -k -H "X-Auth-Token: $token" https://${bmc}/xyz/openbmc_project/inventory/system/chassis/motherboard
\end{verbatim}

\hypertarget{sensors}{%
\subsection{Sensors}\label{sensors}}

The system sensor structure is under the
\texttt{/xyz/openbmc\_project/sensors} hierarchy.

This interface allows monitoring of system attributes like temperature
or altitude, and are represented similar to the inventory, by object
paths under the top-level \texttt{sensors} object name. The path
categorizes the sensor and shows what the sensor represents, but does
not necessarily represent the physical topology of the system.

For example, all temperature sensors are under
\texttt{sensors/temperature}. CPU temperature sensors would be
\texttt{sensors/temperature/cpu{[}n{]}}.

These are some common properties:

\begin{itemize}
\tightlist
\item
  \texttt{Value}: Current value of the sensor
\item
  \texttt{Unit}: Unit of the value and ``Critical'' and ``Warning''
  values
\item
  \texttt{Scale}: The scale of the value and ``Critical'' and
  ``Warning'' values
\item
  \texttt{CriticalHigh} \& \texttt{CriticalLow}: Sensor device
  upper/lower critical threshold bound
\item
  \texttt{CriticalAlarmHigh} \& \texttt{CriticalAlarmLow}: True if the
  sensor has exceeded the critical threshold bound
\item
  \texttt{WarningHigh} \& \texttt{WarningLow}: Sensor device upper/lower
  warning threshold bound
\item
  \texttt{WarningAlarmHigh} \& \texttt{WarningAlarmLow}: True if the
  sensor has exceeded the warning threshold bound
\end{itemize}

A temperature sensor might look like:

\begin{verbatim}
    curl -k -H "X-Auth-Token: $token" https://${bmc}/xyz/openbmc_project/sensors/temperature/ocp_zone
    {
      "data": {
        "CriticalAlarmHigh": false,
        "CriticalAlarmLow": false,
        "CriticalHigh": 65000,
        "CriticalLow": 0,
        "Functional": true,
        "MaxValue": 0,
        "MinValue": 0,
        "Scale": -3,
        "Unit": "xyz.openbmc_project.Sensor.Value.Unit.DegreesC",
        "Value": 34625,
        "WarningAlarmHigh": false,
        "WarningAlarmLow": false,
        "WarningHigh": 63000,
        "WarningLow": 0
      },
      "message": "200 OK",
      "status": "ok"
    }
\end{verbatim}

Note the value of this sensor is 34.625C (34625 * 10\^{}-3).

Unlike IPMI, there are no ``functional'' sensors in OpenBMC; functional
states are represented in the inventory.

To enumerate all sensors in the system:

\begin{verbatim}
curl -k -H "X-Auth-Token: $token" https://${bmc}/xyz/openbmc_project/sensors/enumerate
\end{verbatim}

List properties of one inventory item:

\begin{verbatim}
curl -k -H "X-Auth-Token: $token" https://${bmc}/xyz/openbmc_project/sensors/temperature/outlet
\end{verbatim}

\hypertarget{event-logs}{%
\subsection{Event Logs}\label{event-logs}}

The event log structure is under the
\texttt{/xyz/openbmc\_project/logging/entry} hierarchy. Each event is a
separate object under this structure, referenced by number.

BMC and host firmware on POWER-based servers can report event logs to
the BMC. Typically, these event logs are reported in cases where host
firmware cannot start the OS, or cannot reliably log to the OS.

The properties associated with an event log are as follows:

\begin{itemize}
\tightlist
\item
  \texttt{Message}: The type of event log (e.g.
  ``xyz.openbmc\_project.Inventory.Error.NotPresent'').
\item
  \texttt{Resolved} : Indicates whether the event has been resolved.
\item
  \texttt{Severity}: The level of problem (``Info'', ``Error'', etc.).
\item
  \texttt{Timestamp}: The date of the event log in epoch time.
\item
  \texttt{Associations}: A URI to the failing inventory part.
\end{itemize}

To list all reported event logs:

\begin{verbatim}
    curl -k -H "X-Auth-Token: $token" https://${bmc}/xyz/openbmc_project/logging/entry
    {
      "data": [
        "/xyz/openbmc_project/logging/entry/3",
        "/xyz/openbmc_project/logging/entry/2",
        "/xyz/openbmc_project/logging/entry/1",
        "/xyz/openbmc_project/logging/entry/7",
        "/xyz/openbmc_project/logging/entry/6",
        "/xyz/openbmc_project/logging/entry/5",
        "/xyz/openbmc_project/logging/entry/4"
      ],
      "message": "200 OK",
      "status": "ok"
    }
\end{verbatim}

To read a specific event log:

\begin{verbatim}
    curl -k -H "X-Auth-Token: $token" https://${bmc}/xyz/openbmc_project/logging/entry/1
    {
      "data": {
        "AdditionalData": [
          "_PID=183"
        ],
        "Id": 1,
        "Message": "xyz.openbmc_project.Common.Error.InternalFailure",
        "Purpose": "xyz.openbmc_project.Software.Version.VersionPurpose.BMC",
        "Resolved": false,
        "Severity": "xyz.openbmc_project.Logging.Entry.Level.Error",
        "Timestamp": 1563191362822,
        "Version": "2.8.0-dev-132-gd1c1b74-dirty",
        "associations": []
      },
      "message": "200 OK",
      "status": "ok"
    }
\end{verbatim}

To delete an event log (log 1 in this example), call the \texttt{Delete}
method on the event:

\begin{verbatim}
curl -k -H "X-Auth-Token: $token" -H "Content-Type: application/json" -X POST -d '{"data" : []}' https://${bmc}/xyz/openbmc_project/logging/entry/1/action/Delete
\end{verbatim}

To clear all event logs, call the top-level \texttt{DeleteAll} method:

\begin{verbatim}
curl -k -H "X-Auth-Token: $token" -H "Content-Type: application/json" -X POST -d '{"data" : []}' https://${bmc}/xyz/openbmc_project/logging/action/DeleteAll
\end{verbatim}

\hypertarget{host-boot-options}{%
\subsection{Host Boot Options}\label{host-boot-options}}

With OpenBMC, the Host boot options are stored as D-Bus properties under
the \texttt{control/host0/boot} path. Properties include
\href{https://github.com/openbmc/phosphor-dbus-interfaces/blob/master/yaml/xyz/openbmc_project/Control/Boot/Mode.interface.yaml}{\texttt{BootMode}},
\href{https://github.com/openbmc/phosphor-dbus-interfaces/blob/master/yaml/xyz/openbmc_project/Control/Boot/Source.interface.yaml}{\texttt{BootSource}}
and if the host is based on x86 CPU also
\href{https://github.com/openbmc/phosphor-dbus-interfaces/blob/master/yaml/xyz/openbmc_project/Control/Boot/Type.interface.yaml}{\texttt{BootType}}.

\begin{itemize}
\tightlist
\item
  Set boot mode:
\end{itemize}

\begin{verbatim}
curl -k -H "X-Auth-Token: $token" -H "Content-Type: application/json" -X PUT https://${bmc}/xyz/openbmc_project/control/host0/boot/attr/BootMode -d '{"data": "xyz.openbmc_project.Control.Boot.Mode.Modes.Regular"}'
\end{verbatim}

\begin{itemize}
\tightlist
\item
  Set boot source:
\end{itemize}

\begin{verbatim}
curl -k -H "X-Auth-Token: $token" -H "Content-Type: application/json" -X PUT https://${bmc}/xyz/openbmc_project/control/host0/boot/attr/BootSource -d '{"data": "xyz.openbmc_project.Control.Boot.Source.Sources.Default"}'
\end{verbatim}

\begin{itemize}
\tightlist
\item
  Set boot type (valid only if host is based on the x86 CPU):
\end{itemize}

\begin{verbatim}
curl -k -H "X-Auth-Token: $token" -H "Content-Type: application/json" -X PUT https://${bmc}/xyz/openbmc_project/control/host0/boot/attr/BootType -d '{"data": "xyz.openbmc_project.Control.Boot.Type.Types.EFI"}'
\end{verbatim}

Also there are boolean
\href{https://github.com/openbmc/phosphor-dbus-interfaces/blob/master/yaml/xyz/openbmc_project/Object/Enable.interface.yaml}{\texttt{Enable}}
properties that control if the boot source override is persistent or
one-time, and if the override is enabled or not.

\begin{itemize}
\tightlist
\item
  Set boot override one-time flag:
\end{itemize}

\begin{verbatim}
curl -k -H "X-Auth-Token: $token" -H "Content-Type: application/json" -X PUT https://${bmc}/xyz/openbmc_project/control/host0/boot/one_time/attr/Enabled -d '{"data": "true"}'
\end{verbatim}

\begin{itemize}
\tightlist
\item
  Enable boot override:
\end{itemize}

\begin{verbatim}
curl -k -H "X-Auth-Token: $token" -H "Content-Type: application/json" -X PUT https://${bmc}/xyz/openbmc_project/control/host0/boot/attr/Enabled -d '{"data": "true"}'
\end{verbatim}

\hypertarget{host-state-control}{%
\subsection{Host State Control}\label{host-state-control}}

The host can be controlled through the \texttt{host} object. The object
implements a number of actions including power on and power off. These
correspond to the IPMI \texttt{power\ on} and \texttt{power\ off}
commands.

Assuming you have logged in, the following will power on the host:

\begin{verbatim}
curl -k -H "X-Auth-Token: $token" -H "Content-Type: application/json" -d '{"data": "xyz.openbmc_project.State.Host.Transition.On"}' -X PUT https://${bmc}/xyz/openbmc_project/state/host0/attr/RequestedHostTransition
\end{verbatim}

To power off the host:

\begin{verbatim}
curl -k -H "X-Auth-Token: $token" -H "Content-Type: application/json" -d '{"data": "xyz.openbmc_project.State.Host.Transition.Off"}' -X PUT https://${bmc}/xyz/openbmc_project/state/host0/attr/RequestedHostTransition
\end{verbatim}

To issue a hard power off (accomplished by powering off the chassis):

\begin{verbatim}
curl -k -H "X-Auth-Token: $token" -H "Content-Type: application/json" -X PUT -d '{"data":"xyz.openbmc_project.State.Chassis.Transition.Off"}' https://${bmc}//xyz/openbmc_project/state/chassis0/attr/RequestedPowerTransition
\end{verbatim}

To reboot the host:

\begin{verbatim}
curl -k -H "X-Auth-Token: $token" -H "Content-Type: application/json" -X PUT -d '{"data":"xyz.openbmc_project.State.Host.Transition.Reboot"}' https://${bmc}/xyz/openbmc_project/state/host0/attr/RequestedHostTransition
\end{verbatim}

More information about Host State Management can be found here:
\url{https://github.com/openbmc/phosphor-dbus-interfaces/tree/master/yaml/xyz/openbmc_project/State}

\hypertarget{host-clear-gard}{%
\subsection{Host Clear GARD}\label{host-clear-gard}}

On OpenPOWER systems, the host maintains a record of bad or non-working
components on the GARD partition. This record is referenced by the host
on subsequent boots to determine which parts should be ignored.

The BMC implements a function that simply clears this partition. This
function can be called as follows:

\begin{itemize}
\tightlist
\item
  Method 1: From the BMC command line:
\end{itemize}

\begin{verbatim}
  busctl call org.open_power.Software.Host.Updater \
    /org/open_power/control/gard \
    xyz.openbmc_project.Common.FactoryReset Reset
\end{verbatim}

\begin{itemize}
\tightlist
\item
  Method 2: Using the REST API:
\end{itemize}

\begin{verbatim}
curl -k -H "X-Auth-Token: $token" -H "Content-Type: application/json" -X POST -d '{"data":[]}' https://${bmc}/org/open_power/control/gard/action/Reset
\end{verbatim}

Implementation:
\url{https://github.com/openbmc/openpower-pnor-code-mgmt}

\hypertarget{host-watchdog}{%
\subsection{Host Watchdog}\label{host-watchdog}}

The host watchdog service is responsible for ensuring the host starts
and boots within a reasonable time. On host start, the watchdog is
started and it is expected that the host will ping the watchdog via the
inband interface periodically as it boots. If the host fails to ping the
watchdog within the timeout then the host watchdog will start a systemd
target to go to the quiesce target. System settings will then determine
the recovery behavior from that state, for example, attempting to reboot
the system.

The host watchdog utilizes the generic
\href{https://github.com/openbmc/phosphor-watchdog}{phosphor-watchdog}
repository. The host watchdog service provides 2 files as configuration
options into phosphor-watchdog:

\begin{verbatim}
    /lib/systemd/system/phosphor-watchdog@poweron.service.d/poweron.conf
    /etc/default/obmc/watchdog/poweron
\end{verbatim}

\texttt{poweron.conf} contains the ``Conflicts'' relationships to ensure
the watchdog service is stopped at the correct times. \texttt{poweron}
contains the required information for phosphor-watchdog (more
information on these can be found in the
\href{https://github.com/openbmc/phosphor-watchdog}{phosphor-watchdog}
repository).

The 2 service files involved with the host watchdog are:

\begin{verbatim}
    phosphor-watchdog@poweron.service
    obmc-enable-host-watchdog@0.service
\end{verbatim}

\texttt{phosphor-watchdog@poweron} starts the host watchdog service and
\texttt{obmc-enable-host-watchdog} starts the watchdog timer. Both are
run as a part of the \texttt{obmc-host-startmin@.target}. Service
dependencies ensure the service is started before the enable is called.

The default watchdog timeout can be found within the
\href{https://github.com/openbmc/phosphor-dbus-interfaces/blob/master/yaml/xyz/openbmc_project/State/Watchdog.interface.yaml}{dbus
interface specification} (Interval property).

The host controls the watchdog timeout and enable/disable once it
starts.
